% EGA V, English translation by P. Blass, 1987 (cote 24): archive cote 24, site page 65 (PDF page 66).
% typescript. Draft compared: none.
% AI-transcribed by Claude Opus 5.5 (claude-opus-5-5, adaptive thinking, effort high), Message Batches API, batch ega24b (msgbatch_01RGJJeHdR96U16ZWFuUjtDx), 2026-11-13; not yet checked by eye.
\scanpage{24/65}
\mspage{65}

The proof is essentially trivial in terms of the criteria of
the end of the previous section. Taking an affine model of $P$
containing $f(X)$ we are reduced a) to finding 3 polynomials in the
coordinates $T_1$, $T_2$,....., $T_s$ of degree $\leq 2$, let them be $P$, $Q$, and $R$
such that $Q(t)/P(t)$ and $R(t)/P(t)$ are algebraically independent
over $k$ in $K$ (where $K$ is the function field of $X$ and $t = (t_1$,....,
$t_r)$ is the system of elements of $K$ defined by the $T_i$); in b) we
require also that $P$, $Q$, and $R$ should vanish to order two at least
on $L$ which we can in addition suppose to be defined by the
equations $T_1$,...., $T_s = 0$; finally in c) it is the same but $L$ is
the image of the tangent space to $X$ at $x$ and we allow possibly to
take $P$, $Q$, and $R$ of degree 3 i.e. move margin, the hypothesis that
dim $f(X) \geq 2$ signifies that the transcendence degree of $K(t_1$,
$t_2$,...., $t_r)$ over $k$ is $\geq 2$ i.e. we can find $t_1$, $t_2$ let us say
algebraically independent. In a) we take therefore $P = T_1$, $Q =$
$T_1{}^2$, $R = T_1T_2$, in b) we (analogously) do the same noting that, we
may there choose $t_1$ leading to $T_1$ zero over $L$ due to the fact that
$f(X)$ C $L$ (which implies that there exists an index $i$ between one
and $s$ such that $t_s \neq 0$, so that $t_s$ is not a constant (since $t_s$ is
zero at $x$) therefore $t_s$ is not algebraic over $k$ (\uline{???}) (N.B. we may
suppose $k$ algebraically closed). The case c) follows from b)
except in the case where $f(X)$ is contained in the image of $L$ by $f$
of the tangent space to $X$ at $x$\ \ \ . [If dim $X = 2$ this case is
effectively exceptional (the trace of a quadric surface tangent to
a plane on that plane is in general formed by two intersecting
lines and is therefore not irreducible]. But to treat that case,
in the forms $P$, $Q$, and $R$ made explicit above we may replace
